\documentclass[11pt]{article}
\usepackage{icmlko}
\begin{document}

\title{외우는 용량과 배우는 용량 --- 같은 적합도를 사는데 하나는
집 밖에서 안 쓰인다}
\author{월드모델 실험실}
\date{노트 461 · 2026년 8월 2일}
\maketitle

\begin{abstract}
노트 460 이 이상한 것을 남겼다. 깊이와 \texttt{it} 은 둘 다 ``용량''인데
집 밖에 대해 \textbf{부호가 반대}다 --- 깊이를 늘리면 집 밖이 오르고
\texttt{it} 을 늘리면 내려간다. 노트 $408 \sim 410$ 이 한 가족이라 부른
손잡이들이다.

가설은 하나였다. \emph{깊이는 상호작용을 늘리고 \texttt{it} 은 잔차를 더
깎는데, 잔차를 깎는 것이 그 수집 고유의 무늬를 외우는 일이라면 집 밖에서
손해가 난다.} 그래서 \textbf{학습 행 자신에 대한 적합도}(외운 정도)를
같이 찍었다.

\begin{center}
\begin{tabular}{lrrr}
\hline
설정 & 학습 적합도 & 유보 판 & KR 만화 \\
\hline
깊이 12 · \texttt{it} $120$ & $+0.7188$ & $+0.4795$ & $+0.5612$ \\
깊이 12 · \texttt{it} $220$ & $+0.7604$ & $+0.4827$ & $+0.5580$ \\
깊이 12 · \texttt{it} $400$ & $+0.8038$ & $+0.4806$ & $+0.5491$ \\
깊이 12 · \texttt{it} $700$ & $+0.8426$ & $+0.4767$ & $+0.5382$ \\
\hline
깊이 $6$ · \texttt{it} 220 & $+0.7376$ & $+0.4860$ & $+0.5526$ \\
깊이 $12$ · \texttt{it} 220 & $+0.7604$ & $+0.4827$ & $+0.5580$ \\
깊이 $24$ · \texttt{it} 220 & $+0.7611$ & $+0.4820$ & $+0.5590$ \\
\hline
\end{tabular}
\end{center}

\textbf{갈린다.} \texttt{it} 은 적합도를 $+0.124$ 사면서 집 밖을
$-0.023$ 판다. 깊이는 적합도를 $+0.024$ 밖에 못 사는데 집 밖을
$+0.006$ 산다. \textbf{같은 저울에서 산 적합도인데 하나는 집 밖에서
쓰이고 하나는 안 쓰인다.}
\end{abstract}

\begin{figure}[t]
\centering
\includegraphics[width=\textwidth]{figs/f1.pdf}
\caption{\textbf{가} --- 가로가 외운 정도, 세로가 집 밖 전이. 두 손잡이가
같은 평면에서 \emph{반대 방향}으로 간다. \textbf{나} --- \texttt{it} 은
적합도를 다섯 배 사고 그 대가로 집 밖을 판다.}
\end{figure}

\section{다섯 배 사고 거꾸로 간다}

숫자를 한 줄로 줄이면 이렇다.

\begin{center}
\begin{tabular}{lrrr}
\hline
 & 산 적합도 & 집 밖 변화 & 적합도 한 단위당 \\
\hline
\texttt{it} $120 \to 700$ & $+0.1238$ & $-0.0230$ & $\mathbf{-0.19}$ \\
깊이 $6 \to 24$ & $+0.0235$ & $+0.0064$ & $\mathbf{+0.27}$ \\
\hline
\end{tabular}
\end{center}

\textbf{적합도 한 단위를 살 때 \texttt{it} 은 집 밖을 $0.19$ 잃고 깊이는
$0.27$ 번다.} 부호가 반대이고 크기도 비슷하다.

그리고 모양이 다르다. \textbf{깊이의 적합도는 포화한다} --- $12 \to 24$
에서 $+0.0007$ 밖에 안 는다(그래도 집 밖은 $+0.0010$ 오른다).
\textbf{\texttt{it} 의 적합도는 안 포화한다} --- $700$ 까지 계속 올라
$0.84$ 에 이르고, 그동안 집 밖은 계속 내려간다.

\emph{깊이는 살 것을 다 사면 멈추고, \texttt{it} 은 살 것이 없어도 계속
산다.} 없는 것을 계속 사면 그것이 외우기다.

\section{그래서 ``용량''이라는 말을 나눈다}

노트 $408 \sim 410$ 은 깊이 · lr · 자루섞기가 \textbf{같은 용량 가족}이라
했다 --- 도메인이 그 셋에 매기는 선호 순서가 같았기 때문이다(정체는
``F18$-$선형 격차'', rho $+0.773$). 그 관찰은 여전히 맞다. 하지만
\textbf{선호 순서가 같다고 같은 것을 사는 건 아니다.}

\begin{quote}
\textbf{배우는 용량}(깊이) --- 축들 사이의 \emph{상호작용}을 표현한다.
적합도를 조금 사고 포화하며, 그 조금이 \textbf{집 밖에서도 쓰인다}. \\
\textbf{외우는 용량}(\texttt{it}) --- 남은 \emph{잔차}를 깎는다. 적합도를
많이 사고 포화하지 않으며, 그 대부분이 \textbf{그 수집 안에서만 쓰인다}.
\end{quote}

노트 459 가 ``환율이 손잡이마다 다르다''고 적었는데, 이 노트가 그
\emph{왜}를 준다. \textbf{환율이 다른 게 아니라 사는 물건이 다르다.}

\section{규약에 닿는 곳}

노트 460 이 조건 ④(환율 $\ge 1$)를 넣었다. 이 노트는 그 조항을
안 바꾸지만 \textbf{쓰는 법을 하나 더한다}.

\begin{quote}
\textbf{조항} --- 용량을 바꾸는 변경은 \textbf{학습 적합도도 같이 적는다.}
적합도를 많이 사면서 집 밖을 잃는 변경은 \emph{환율이 우연히 $1$ 을
넘더라도} 외우기를 사고 있는 것이니 근거를 따로 대야 한다.
\end{quote}

\section*{남은 것}

\begin{itemize}
\item \textbf{잔차를 깎는 것이 정말 ``수집 고유''인가.} 지금은 간접
증거다 --- 적합도는 오르는데 집 밖이 내려간다. 직접 보려면 \texttt{it}
이 늘 때 나무가 \emph{원핫에 얼마나 기대는지}를 재면 된다.
\item \textbf{깊이의 포화가 곧 봉우리다.} $12 \to 24$ 에서 적합도가
$+0.0007$ 인데 집 밖이 $+0.0010$ 이다 --- 노트 458 이 씨앗 열로 기각한
자리다. 적합도 포화가 \textbf{봉우리를 미리 알려 주는 값싼 신호}일 수
있다. 다른 손잡이에서도 되나.
\item \textbf{lr 은 어느 쪽인가.} 노트 459 는 lr 을 기각만 했지
\emph{무엇을 사는지}는 안 봤다.
\end{itemize}

\end{document}
